\section{Architecture of the Multi-Agent System}\label{sec:3}

JobMatch organizes hiring support around three collaborating agents and a shared representation layer.
Figure~\ref{fig:1} shows the overall layout; Figures~\ref{fig:2} and~\ref{fig:3} summarize end-user workflows.
Three design principles guide the layout: \textbf{separation of ownership} (each market side controls its own records), \textbf{read-only matchmaking} (scoring never edits profiles or acts on behalf of users), and \textbf{human-in-the-loop decisions} (save, apply, contact, and shortlist remain explicit user actions).
Concrete libraries, APIs, and parameters appear in \S\ref{sec:4}; score formulas and evaluation metrics in \S\ref{sec:5}.

\begin{JFigure}{../figures/Fig1.pdf}
\figcap{1}{Multi-agent architecture of the proposed JobMatch recruitment system. The Candidate and Employer agents each own an ingestion pipeline (parse, embed, store, profile state); the read-only Matchmaking agent queries both vector stores, scores pairings, and returns ranked recommendations to role portals. A shared communication layer carries events, versioned snapshots, and cache-invalidation signals between agents.}
\end{JFigure}

\subsection{Candidate/Client Agent}\label{sec:3-candidate}

The \textbf{candidate-side agent} (also referred to as the client-side agent in our implementation) represents the job seeker.
Its responsibility is to turn resumes and profile edits into a maintained \emph{candidate state} that the rest of the system can trust.

When a user uploads or revises a resume, this agent assists with parsing, normalizes skills into a shared vocabulary, and registers a versioned snapshot of the person's skills, experience, compensation expectations, and remote preferences.
It also maintains the editable portal profile the user sees in forms.
The agent publishes a profile-updated signal whenever a new snapshot is ready so that stale match lists can be refreshed.

The candidate agent does \emph{not} own job postings, write into employer records, or decide which applications are sent.
Those boundaries keep accountability clear: the seeker's representative manages seeker data only.

\subsection{Employer Agent}\label{sec:3-employer}

The \textbf{employer-side agent} represents the hiring organization.
It accepts job descriptions and requirement updates, structures them into a maintained \emph{job state}, and keeps that state aligned with what recruiters edit in the employer portal.

For each posting, the agent registers a snapshot with required and preferred skills, experience bounds, compensation band, remote policy, and descriptive text suitable for comparison against candidates.
Like its counterpart on the candidate side, it emits a job-updated signal when a snapshot changes so match rankings can be invalidated without cross-agent writes.

The employer agent does \emph{not} manage candidate profiles or override match scores.
Shortlisting and applicant review remain human decisions supported by ranked suggestions.

\leadpara{Optional external job feed.}
When enabled at deployment time, the Employer agent can also ingest paginated job listings from an external provider API.
Each raw record is normalized to the same canonical job schema used for portal postings, deduplicated by stable identifier, and written to a local snapshot file for restart recovery.
A sync operation replaces the in-memory employer corpus and re-embeds job snapshots; the frozen \texttt{jobs.json} demo seed remains the default for offline evaluation reported in \S\ref{sec:6}.
The Matchmaking agent consumes the refreshed snapshots read-only; it does not call the external provider directly.

\subsection{Matchmaking Agent}\label{sec:3-matchmaking}

The \textbf{matchmaking agent} is a neutral broker between the two market sides.
Given a candidate snapshot and a set of job snapshots---or the reverse direction for employer-initiated search---it scores each feasible pairing, attaches structured explanations, and returns a ranked list.

It reads snapshots by identifier only; it never mutates candidate or job stores.
It may fuse multiple retrieval signals and optional reranking stages when configured, but the default portal path favors a fixed composite strategy whose components users can inspect (\S\ref{sec:4}, \S\ref{sec:5}).
When profile or job events indicate that inputs changed, the matchmaking layer drops cached rankings for the affected entity and recomputes on the next request.

The matchmaking agent suggests; it does not hire, reject, or apply on anyone's behalf.

\subsection{Agent Communication and State Sharing}\label{sec:3-communication}

JobMatch relies on exactly three cooperating agents: the \textbf{Candidate/Client Agent}, the \textbf{Employer Agent}, and the \textbf{Matchmaking Agent} (Figure~\ref{fig:1}).
Sections~\ref{sec:3-candidate}--\ref{sec:3-matchmaking} describe each agent's role; this subsection explains how they exchange information, what each one is allowed to keep private, and how user actions propagate through the system without collapsing ownership boundaries.

\leadpara{What each agent owns.}
The Candidate/Client Agent owns everything derived from the job seeker's side of the market: the resume-derived profile, stated preferences (salary, remote work, location constraints), normalized skills, and the embedding-backed \emph{candidate state} used for comparison.
The Employer Agent owns the parallel objects on the hiring side: the job-description-derived posting, structured requirements, compensation and experience bounds, and the \emph{job state} with its embedding.
The Matchmaking Agent owns neither profile store.
It owns the retrieval-and-ranking workflow---how candidate and job snapshots are paired, scored, explained, and ordered---and the transient recommendation output returned to a portal session.
It may cache recent rankings for responsiveness, but it does not author candidate or job records.

\leadpara{State each agent maintains.}
On the candidate and employer sides, each agent maintains two related layers.
A \textbf{profile} is the editable portal record: form fields the user can revise at any time.
A \textbf{snapshot} is a versioned, match-ready view frozen at the moment the user confirms a save---normalized skills, constraint values, document identity, and a comparable text representation with its embedding.
Only snapshots enter ranking; profiles exist so people can edit without immediately rewriting history.
The Matchmaking Agent maintains session-scoped ranking state (cached lists, last-computed scores) and structured explanation payloads attached to each recommendation.
It does not maintain a third copy of full resumes or job descriptions beyond what it needs to score a request.

\leadpara{What data is shared.}
Agents share the minimum information required for fair comparison and timely refresh.
\textbf{Snapshots and identifiers} cross the boundary to the Matchmaking Agent when a user requests matches: a candidate identifier plus the current candidate snapshot, and a set of job identifiers with their snapshots (or the reverse direction for employer search).
\textbf{Events} broadcast that a snapshot changed---for example, after a confirmed profile save or job edit---so the Matchmaking Agent can discard stale rankings.
\textbf{Canonical skill vocabulary} is shared implicitly: both owning agents normalize skills to the same alias table so that overlap and constraint checks refer to the same terms.
\textbf{Recommendation results} flow outward as ranked lists with explanations; they are shared with the UI layer, not written back into profile stores.

\leadpara{What should not be shared unnecessarily.}
The design deliberately withholds data that would blur accountability or expose one party without cause.
The Candidate Agent does not receive employer-only draft postings or internal recruiter notes unless the employer publishes a job snapshot.
The Employer Agent does not receive another candidate's full profile when an employer views a shortlist---only the fields appropriate to that match request.
The Matchmaking Agent never receives portal credentials, authentication secrets, or raw feedback logs as inputs to ranking; save, dismiss, and apply actions are recorded separately and do not silently rewrite embeddings.
Neither owning agent calls the other's write path: all cross-side access for matching is read-only through the Matchmaking Agent.

\leadpara{How vector stores are updated.}
Each owning agent updates only its own embedding index when a confirmed snapshot changes.
After parsing and normalization, the agent computes a fresh vector for the canonical document text and upserts it into the candidate or job vector store.
The opposite side's store is untouched.
An update event follows the write so downstream components know prior similarity searches may be outdated.
The Matchmaking Agent never performs upserts; it queries both stores when retrieving candidates for a job or jobs for a candidate.
Implementation details of the store backends appear in \S\ref{sec:4}.

\leadpara{How the Matchmaking Agent reads from both sides.}
A match request supplies direction (candidate seeking jobs, or employer seeking candidates), the authenticated user's snapshot identifier, and retrieval scope (top-$K$, strategy).
The agent loads the query-side snapshot from the owning agent's store, retrieves candidate vectors or job vectors from the respective indexes, scores each feasible pair using the configured strategy, attaches explanations, and sorts the results.
Reads are symmetric: the same pipeline runs in reverse when an employer searches for candidates.
If either store lacks a fresh embedding, the owning agent must register a snapshot first; the Matchmaking Agent does not fabricate missing state.

\leadpara{How UI actions trigger agent updates.}
Portal interactions map cleanly to owning agents and to the Matchmaking layer.
When a candidate uploads or edits a resume and confirms the form, the Candidate Agent updates profile and snapshot state, refreshes the candidate embedding, and emits a profile-updated event.
When an employer posts or revises a job description and saves, the Employer Agent performs the parallel job-side update.
Requesting the matches page---or an employer opening reverse search for a posting---invokes the Matchmaking Agent with the current snapshots; no ranking runs on half-edited forms until the user confirms.
Secondary actions such as save, apply, shortlist, or dismiss are recorded as user activity; they do not automatically change embeddings unless the user later edits and re-saves a profile or job.
Administrators can inspect recent events for debugging but do not override ownership through the UI.

\leadpara{Human-in-the-loop by design.}
This communication model keeps people in control of outcomes that matter.
Agents prepare, normalize, and suggest; they do not substitute for a hiring decision.
Ranked lists arrive with reasons attached, and every consequential action---saving a job, submitting an application, contacting a candidate---requires an explicit portal choice.
Because snapshots change only after user confirmation, a person can review parsed fields, correct errors, and understand what the system will compare before any recommendation is computed.
Separation of ownership plus read-only matchmaking makes it clear who changed what: candidates govern candidate state, employers govern job state, and the Matchmaking Agent governs only how those states are compared at a given moment.

\subsection{Overall Workflow}\label{sec:3-workflow}

Figure~\ref{fig:2} traces the candidate path; Figure~\ref{fig:3} traces the employer path.
In both portal workflows, document ingestion is assistive: extracted fields populate review forms, users confirm or correct them, and only then does the owning agent register a snapshot and notify the matchmaking layer.

\subsubsection*{Detailed Workflow}\label{workflow:matching}

\small
\noindent\textbf{Inputs:}
Candidate resumes or CVs; job descriptions; candidate preferences when available; employer constraints when available.

\noindent\textbf{Outputs:}
Ranked candidate--job matches; matching scores; structured explanations and matched-skill summaries.

\begin{enumerate}[label=\roman*., leftmargin=1.6em, itemsep=0.12em, topsep=0.2em]
  \item \textbf{Candidate Agent} receives a resume or CV from the candidate portal.
  \item \textbf{Candidate Agent} preprocesses and parses the document into assistive form fields.
  \item \textbf{Candidate Agent} extracts skills, experience, education, and stated preferences; the user confirms or edits the result.
  \item \textbf{Candidate Agent} generates a candidate embedding from the confirmed canonical text.
  \item \textbf{Candidate Agent} stores the candidate profile, versioned snapshot, and vector representation; emit a profile-updated event.
  \item \textbf{Employer Agent} receives a job description from the employer portal.
  \item \textbf{Employer Agent} preprocesses and parses the job description.
  \item \textbf{Employer Agent} extracts role title, required skills, experience bounds, compensation band, location or remote policy, and hiring constraints; the user confirms or edits the result.
  \item \textbf{Employer Agent} generates a job embedding from the confirmed canonical text.
  \item \textbf{Employer Agent} stores the job profile, versioned snapshot, and vector representation; emit a job-updated event.
  \item \textbf{Matchmaking Agent} retrieves candidate and job representations from both vector stores (read-only).
  \item \textbf{Matchmaking Agent} computes semantic similarity and optional constraint-aware component scores (formulas in Section~\ref{sec:5}).
  \item \textbf{Matchmaking Agent} ranks candidate--job pairs by the configured strategy (portal default: composite score).
  \item \textbf{Matchmaking Agent} generates an explanation and matched-skill summary for each ranked pair.
  \item \textbf{UI/Application Layer} displays ranked recommendations; the user explicitly chooses save, apply, shortlist, or contact.
\end{enumerate}
\normalsize

\algref{alg:matching} summarizes the end-to-end multi-agent matching process at the architectural level; score definitions appear in \S\ref{sec:5} and implementation paths in \S\ref{sec:4}.

\begin{JFigure}{../figures/Fig2.pdf}
\figcap{2}{Candidate workflow: resume upload, assisted parsing, profile confirmation, snapshot registration, and ranked job discovery with explainable results.}
\end{JFigure}

\begin{JFigure}{../figures/Fig3.pdf}
\figcap{3}{Employer workflow: job description ingestion, posting confirmation, reverse candidate matching, and applicant review with structured explanations.}
\end{JFigure}

Candidates browse ranked jobs with explanation panels, then explicitly save or apply.
Employers browse ranked candidates for a posting, then shortlist or contact through portal actions.
Administrators inspect agent health and recent events in a separate console.
Separate web experiences for candidates, employers, and administrators sit above the gateway; routing, endpoints, and default request parameters are specified in \S\ref{sec:4}.
